\section{High Level Overview}
This project builds a scalable end-to-end pipeline for detecting vegetation stress preceding Mountain Pine Beetle outbreaks in Boulder County, Colorado. Satellite imagery from three independent datasets is ingested, indexed, transformed, and analyzed using a combination of AWS S3, PostgreSQL/PostGIS, and Apache PySpark. The pipeline is designed so that the full archive is always query able, but only the data relevant to a specific spatiotemporal window is ever loaded into memory.
\begin{itemize}
    \item \textbf{Data Sources and Study Context} --- Data collection using Google Earth Engine and ESA World Cover Datasets
    Gemini said
    \item \textbf{Data Management and Indexing} --- Satellite imagery exported from Google Earth Engine is stored in AWS S3 as Cloud Optimized GeoTIFFs with LZW compression for efficient byte-range reads, with metadata and spatial geometries indexed in a PostgreSQL/PostGIS database for optimized retrieval.
    \item \textbf{Data Transformation Pipeline} --- Matching scene URLs are retrieved from PostGIS, band arrays are streamed directly from S3 COGs, forest-masked at the NumPy level, and aggregated into monthly NDVI/NDMI composites in PySpark.
    \item \textbf{Query-Triggered Anomaly Detection} --- Per-pixel Z-scores are computed against a historical monthly baseline, and pixel-months falling below a threshold of $-1.5$ are flagged as vegetation stress anomalies and written to S3 as Parquet.
    \item \textbf{Evaluation Design} --- Three pipeline runs across serial S3 scan, concurrent S3 scan, and PostGIS GIST-indexed query modes quantify the performance gain from spatial indexing over unfiltered S3 enumeration.
\end{itemize}